\section{Term sheet detail for the six notes}\label{app:terms}

Each note is identified below by CUSIP and by the EDGAR accession number of the
priced supplement it was read from. All six documents are preserved in full in
the project repository, and all six remain retrievable from EDGAR under these
accession numbers. The economic terms appear in Table~\ref{tab:terms}; what
follows is the information needed to locate the source and to see which
structural variant each note is.

\begin{table}[htbp]
\centering
\caption{Filing identifiers and structural variant}
\label{tab:filings}
\small
\begin{tabular}{l l l l l}
\toprule
Note & CUSIP & Accession number & Coupon frequency & Redemption \\
\midrule
N1 & 46661PD38 & 0001918704-26-031125 & monthly   & issuer's option \\
N2 & 61781LWY2 & 0001839882-26-050390 & monthly   & rational-exercise rule \\
N3 & 61781LM23 & 0001839882-26-050393 & monthly   & rational-exercise rule \\
N4 & 40058M7F9 & 0001193125-26-419171 & monthly   & issuer's option \\
N5 & 09712CWN9 & 0001918704-26-030916 & quarterly & automatic at 93.25\% \\
N6 & 46661PL96 & 0001918704-26-031101 & monthly   & issuer's option \\
\bottomrule
\end{tabular}
\notetext{\textit{Notes:} Product names as printed on the covers are
\textit{Callable Contingent Interest Notes} (N1, N6), \textit{Callable
Contingent Income Securities} (N2, N3), \textit{Callable Contingent Coupon
Index-Linked Notes} (N4) and \textit{Contingent Income Auto-Callable Yield
Notes} (N5). Each note is issued by a finance subsidiary and guaranteed by its
parent. The three notes redeemable at the issuer's option become redeemable on
scheduled payment dates after an initial non-call period: from the twelfth
payment date on N1 and from the third on N4 and N6, in each case excluding the
final payment. The two governed by the rational-exercise rule are redeemable on
every stated redemption date, 48 of them on N2 and 30 on N3.}
\end{table}

One structural difference matters for the simulation and is not visible in
Table~\ref{tab:terms}. Five of the notes observe the downside trigger only at
maturity, so a breach is a European condition on the terminal level. The
auto-callable note observes it on any trading day, which is an American
knock-in and is handled in the simulation by a per-path indicator updated at
every daily step rather than at the observation dates alone.
